\section{Local action prices for the stationary construction}
\label{sec:stationary-action-prices}
We verify the localized derivative input of
Section~\ref{sec:uniform-stationary-chart} for the actual unmasked
block action. The action is a sum over sites. Its positive-order
coordinate derivatives have uniform local bounds, although its
scalar value need not have a volume-independent bound.

Fix the block side $L$, the patch, the direct clip, and the kernel
parameters, including its positive frame and predictor cutoffs and
positive penalty coefficient. Fix also an order $n\ge3$ and an
exponential path width $\mu>0$. The incoming path data have row, cubic and quartic
budgets at most $C_0$ at width $\mu$; the loop budget is at most
$\Lambda$. These budgets are the sums of the absolute coefficients
with their specified exponential length weights. All constants
below may depend on this fixed data and on the group, but not on
the auxiliary torus, its forest, or the tag choices.

\subsection{From a finite read set to a localized derivative}
A function based at a site $v$ has read radius $R$ if changing
coordinates anchored farther than $R$ from $v$ leaves the function
unchanged. A derivative containing any one such coordinate is zero:
restrict to a line varying that coordinate, differentiate the
constant function, and commute derivatives. If derivatives through
order $n$ have operator norm at most $A$, then, for $q\le n$,
\begin{equation}\label{eq:stationary-read-price}
 |\partial_{i_1}\cdots\partial_{i_q}u|
 \le Ae^{anR}\exp\!\left(-a\sum_{h=1}^q
                               \operatorname{dist}(v,a_{i_h})\right).
\end{equation}
If every tested coordinate lies in the read set, the sum of their
distances is at most $nR$; otherwise the left side vanishes.
This also holds for vector-valued expressions with norms in place
of absolute values.

For a path atom of total length $t$, the fine forest chart and a
signed motion enlarge its read radius to at most
$2t+1+L$. The path itself travels at most $2t$ in fine units;
reversing an oriented bond can move its initial point by one,
and the coarse replica coordinate is anchored at most $L$ from
the fine bond. This estimate remains valid for winding paths:
periodic distance is bounded by lifted path length.

Choose
\[
                         a=\frac\mu{4(n+1)(1+L)}.
\]
The elementary inequality
\begin{equation}\label{eq:stationary-support-payment}
         e^{\mu t/2}e^{an(2t+1+L)}\le e^{\mu t}e^\mu
\end{equation}
follows by comparing exponents: the second exponent is at most
$\mu t/2+\mu$.
The holomorphic path bounds at width $\mu/2$ give a derivative
price $A_n e^{\mu t/2}$ for a slot, and
$A_{n,k}\operatorname{magnitude}(q)e^{\mu t/2}$ for a degree-$k$
atom $q$. To recall why the constants are uniform, each matrix
transport is a finite product; on the fixed complex neighborhood
its norm is bounded exponentially in path length. The Cauchy
estimate on a fixed smaller neighborhood bounds derivatives by
that exponential times a constant depending only on $n$ and the
neighborhood. Composing with the fine exponential chart uses its
uniform link derivatives and the finite partition formula. The
fixed color projections and multilinear atom contractions contribute
only the constants $A_n,A_{n,k}$.
Combining this estimate with
\eqref{eq:stationary-read-price} and
\eqref{eq:stationary-support-payment} pays the localized derivative
by the original width-$\mu$ coefficient budget.

For an infinite row or atom family, the paid coefficient sequence
is summable. It bounds the values and derivatives through order
$n$ on the common coordinate ball. Uniform convergence and the
fundamental theorem of calculus along coordinate segments justify
termwise differentiation, inductively in the derivative order.
The underlying holomorphic series converges normally on a smaller
complex neighborhood, so it is smooth to all orders there.
Consequently a row based at $v$ has localized derivative cost
$C_0A_n$, and a cubic or quartic family has cost $C_0A_{n,3}$ or
$C_0A_{n,4}$. A loop family has cost $\Lambda A_{n,2}$.
The latter dependence is linear in $\Lambda$.

\subsection{Summing sites without a volume factor}
Suppose a site term $u_v$ has derivatives bounded by
\[
 |\partial_{i_1}\cdots\partial_{i_q}u_v|
             \le B\exp\!\left(-a\sum_h\operatorname{dist}(v,a_{i_h})\right).
\]
Fix a tested coordinate $i_t$ and set
$s=\sum_h\operatorname{dist}(v,a_{i_h})$. Then
$\operatorname{dist}(v,a_{i_t})\le s$ and
\[
 \sum_h\operatorname{dist}(a_{i_t},a_{i_h})
                                         \le(q+1)s\le(n+1)s.
\]
Splitting the exponent in two gives
\[
 e^{-as}\le e^{-(a/2)\operatorname{dist}(v,a_{i_t})}
     \exp\!\left(-\frac a{2(n+1)}
                     \sum_h\operatorname{dist}(a_{i_t},a_{i_h})\right).
\]
Summing $v$ costs only the distance-shell sum at width $a/2$.
If at most $h$ terms have the same base site, multiply this cost
by $h$. We have obtained the bound
\eqref{eq:stationary-scatter-input}, with width $a/[2(n+1)]$,
for every positive derivative of the site sum. The argument needs
a tested coordinate as root, so it makes no assertion at order
zero. Averaging over signed motions preserves the cost: each
motion permutes sites and the averaging weights sum to one.

The classical site density has three types of terms. The direct
plaquette term has a fixed read radius and derivative bounds from
its four unitary factors and the smooth group penalty. The
squared-row term is $\|\ell(U)\|^2/2$. Its row value and derivatives
are bounded by the row budget just proved; the finite product rule
bounds its derivatives by a constant depending on $C_0$. Its
localized bounds are preserved by that product rule, since the
exponential factors for the split derivative slots multiply to the
factor for all slots. Finally the cubic and quartic terms have the
atom bounds above. Adding these terms and summing sites proves
the required localized derivative estimate for the fine classical
action. The classical price depends on $C_0$ and the direct clip,
and does not depend on the loop budget $\Lambda$.

\subsection{The prediction penalty on the same ball}
The minimizing frame and predictor of
Section~\ref{sec:predictor-branch} are smooth near the identity.
The finite-read transplantation argument of
Section~\ref{sec:mean-action-jets} makes their fixed-order
neighborhood and derivative bounds uniform over periods. We must
also verify that all tags select that same branch there.

On a coordinate ball of radius $r$, each fine bond chord is at
most $C_1r$, by integrating its first derivative from the identity.
Every path used by a frame is of bounded length for the fixed
patch, so its transport chord is at most $C_2r$. Shrink $r$ until
the strict frame test holds. The uniqueness of the normalized
minimizing frame identifies all frame tags there. Its coordinates
and first derivatives are uniformly bounded on a smaller ball.
A connector contains the frame at each of its two endpoints and
a straight fine path. Its chord is therefore at most $C_3r$,
with a common constant $C_3>0$.

Let $Y_i$ be these connectors and $w_i\ge0$ their weights,
$\sum_iw_i=1$. Comparing the prediction cost with the identity gives
\[
       \min_u\sum_iw_i\operatorname{price}(Y_i u^{-1})
          \le\tfrac12(C_3r)^2.
\]
Choose $r<\sqrt{\zeta}/C_3$, where $\zeta>0$ is the predictor
cutoff. This inequality is then strictly below the cutoff. Also
\[
       \left\|\sum_iw_i\operatorname{mat}(Y_i)-1\right\|
                       \le c_GC_3r,
\]
where $c_G$ converts chord to matrix norm. Shrink $r$ further so
this lies in the fixed nearest-point tube. The selected predictor
is the smooth nearest-point projection of this weighted matrix
average. Hence every predictor tag agrees with this same smooth
expression throughout the ball. Product and composition rules
bound its derivatives uniformly, using the same actual endpoint
frames in both occurrences.

For each coarse bond, the smoothing penalty is the group price of
the coarse field times the inverse predicted field. Its value and
its derivatives through order $n$ have a common bound on this ball.
Its read set is the finite union of the patch paths at the two
endpoints and the coarse bond coordinate. Apply
\eqref{eq:stationary-read-price} with that fixed radius, then the
site-summation argument. The embedded coarse sites are distinct;
there are four coarse bond directions at each such site. Thus this
sum has the same localized derivative bound, with multiplicity
four, multiplied by the fixed kernel penalty coefficient.

The full tested action is the sum of this smoothing penalty and
the fine classical action. Taking the smaller common radius and
adding their costs proves
\eqref{eq:stationary-scatter-input} for the actual action, with all
tags equal there. At the identity the curvature slots, plaquette
cost and prediction discrepancy vanish; their first derivatives
in the full action vanish as well, since the direct penalty is
stationary and every other term has degree at least two in these
vanishing slots. Thus $f(0)=0$ and $D_yf(0)=0$.
Together with the tree-gauge Hessian gap, this supplies all inputs
to the uniform stationary construction. The argument fixes the
block and path budgets first; it does not assert a new estimate
for the evolution of those budgets along a history.
